% GENERATED FILE. Edit canonical lessons, then run npm run generate:books.
% canonical-source-hash: fnv1a64:3bc959ab4944da1d
% canonical-lessons: PA-C315-daku, PA-C315-hamla, PA-C315-bhannna, PA-C315-gahina, PA-C315-jan, PA-R315-review-309-315

\chapter{Bandits on the Road}
\label{ch:pa-a2-bandits-on-the-road}
\hlchaptermodality{315}

\begin{chapteropening}
\textbf{By the end of this chapter:} \emph{I can say what happened when someone was robbed.}

Complete the last lesson of chapter 315: I can say what happened when someone was robbed.
\end{chapteropening}

\section[\texorpdfstring{\pa{ਡਾਕੂ} (ḍākū)}{ਡਾਕੂ (ḍākū)}]{\pa{ਡਾਕੂ} (ḍākū) — a bandit}

\label{lesson:PA-C315-daku}



\begin{quote}
Before the new one: say the Punjabi for a son's son, then the Punjabi for a daughter's son.
\end{quote}

\subsection*{You'll want to know: \pa{ਡਾਕੂ}}
\textbf{\pa{ਡਾਕੂ}} — \emph{ḍākū} — \textquotedblleft{}a bandit\textquotedblright{}.

\begin{quote}
\pa{ਪਿਛਲੇ} \pa{ਸਾਲ} \pa{ਸਾਡੇ} \pa{ਕਸਬੇ} \pa{ਵਿੱਚ} \pa{ਡਾਕੂ} \pa{ਆਏ।}
\end{quote}

\begin{quote}
\emph{pichhle sāl sāḍe kasbe vich ḍākū āe.}
\end{quote}

\begin{quote}
Last year bandits came to our small town.
\end{quote}

A masculine noun.

\begin{cousinweb}
\textbf{\pa{ਡਾਕੂ}} (\emph{ḍākū}) is inherited from a reconstructed \emph{*ḍākka}, \textquotedblleft{}a robber\textquotedblright{}; CDIAL lists Punjabi \emph{ḍākū}, \textquotedblleft{}robber\textquotedblright{}, under it (CDIAL 5543).
\end{cousinweb}

\begin{grammarlens}[title={What you've built}]
One more word for saying what happened.
\end{grammarlens}

\subsection*{Your turn}
\begin{itemize}
\raggedright
  \item \emph{Say it:} \emph{ḍākū}
  \item \emph{Say it:} \emph{ḍākū}, once more
  \item \emph{Say it:} \emph{pichhle sāl sāḍe kasbe vich ḍākū āe.} — Last year bandits came to our small town.
  \item \emph{Recall:} say the Punjabi for a son's son, then the Punjabi for a daughter's son, then say \emph{ḍākū} again
\end{itemize}

\begin{tcolorbox}[breakable,colback=teal!4,colframe=teal!35!black,title={Before you move on}]
What does \pa{ਡਾਕੂ} mean? (\textquotedblleft{}A bandit\textquotedblright{}.) Say it once more.
\end{tcolorbox}
\section[\texorpdfstring{\pa{ਹਮਲਾ} (hamlā)}{ਹਮਲਾ (hamlā)}]{\pa{ਹਮਲਾ} (hamlā) — an attack}

\label{lesson:PA-C315-hamla}



\begin{quote}
Before the new one: say the Punjabi for a daughter's son, then the Punjabi for a bandit.
\end{quote}

\subsection*{You'll want to know: \pa{ਹਮਲਾ}}
\textbf{\pa{ਹਮਲਾ}} — \emph{hamlā} — \textquotedblleft{}an attack\textquotedblright{}.

\begin{quote}
\pa{ਰਾਤ} \pa{ਨੂੰ} \pa{ਬੱਸ} \pa{ਉੱਤੇ} \pa{ਹਮਲਾ} \pa{ਹੋਇਆ।}
\end{quote}

\begin{quote}
\emph{rāt nū̃ bass utte hamlā hoiā.}
\end{quote}

\begin{quote}
At night there was an attack on the bus.
\end{quote}

A masculine noun.

\begin{cousinweb}
\textbf{\pa{ਹਮਲਾ}} (\emph{hamlā}) is a loan, usually traced through Persian to Arabic \emph{ḥamla}, \textquotedblleft{}an attack, a charge\textquotedblright{}.
\end{cousinweb}

\begin{grammarlens}[title={What you've built}]
One more word for saying what happened.
\end{grammarlens}

\subsection*{Your turn}
\begin{itemize}
\raggedright
  \item \emph{Say it:} \emph{hamlā}
  \item \emph{Say it:} \emph{hamlā}, once more
  \item \emph{Say it:} \emph{rāt nū̃ bass utte hamlā hoiā.} — At night there was an attack on the bus.
  \item \emph{Recall:} say the Punjabi for a daughter's son, then the Punjabi for a bandit, then say \emph{hamlā} again
  \item \emph{Write it:} copy \textbf{\pa{ਹਮਲਾ}} once, checking the whole word
\end{itemize}

\begin{tcolorbox}[breakable,colback=teal!4,colframe=teal!35!black,title={Before you move on}]
What does \pa{ਹਮਲਾ} mean? (\textquotedblleft{}An attack\textquotedblright{}.) Say it once more.
\end{tcolorbox}
\section[\texorpdfstring{\pa{ਭੰਨਣਾ} (bhannṇā)}{ਭੰਨਣਾ (bhannṇā)}]{\pa{ਭੰਨਣਾ} (bhannṇā) — to smash}

\label{lesson:PA-C315-bhannna}



\begin{quote}
Before the new one: say the Punjabi for a bandit, then the Punjabi for an attack.
\end{quote}

\subsection*{You'll want to know: \pa{ਭੰਨਣਾ}}
\textbf{\pa{ਭੰਨਣਾ}} — \emph{bhannṇā} — \textquotedblleft{}to smash\textquotedblright{}.

\begin{quote}
\pa{ਅਮਨ} \pa{ਨੇ} \pa{ਖਿੜਕੀ} \pa{ਦਾ} \pa{ਸ਼ੀਸ਼ਾ} \pa{ਭੰਨਿਆ।}
\end{quote}

\begin{quote}
\emph{Aman ne khiṛkī dā shīshā bhanniā.}
\end{quote}

\begin{quote}
Aman smashed the windowpane.
\end{quote}

Past (I …): \pa{ਮੈਂ} \pa{ਪੁਰਾਣਾ} \pa{ਘੜਾ} \pa{ਭੰਨਿਆ}, \emph{maiṁ purāṇā ghaṛā bhanniā}.

\begin{cousinweb}
\textbf{\pa{ਭੰਨਣਾ}} (\emph{bhannṇā}) is inherited from Sanskrit \emph{bhajyate}, \textquotedblleft{}is broken\textquotedblright{}; CDIAL lists Punjabi \emph{bhanṇā}, \textquotedblleft{}to break\textquotedblright{}, under it (CDIAL 9361).
\end{cousinweb}

\begin{grammarlens}[title={What you've built}]
One more word for saying what happened.
\end{grammarlens}

\subsection*{Your turn}
\begin{itemize}
\raggedright
  \item \emph{Say it:} \emph{bhannṇā}
  \item \emph{Say it:} \emph{bhannṇā}, once more
  \item \emph{Say it:} \emph{Aman ne khiṛkī dā shīshā bhanniā.} — Aman smashed the windowpane.
  \item \emph{Recall:} say the Punjabi for a bandit, then the Punjabi for an attack, then say \emph{bhannṇā} again
  \item \emph{Write it:} copy \textbf{\pa{ਭੰਨਣਾ}} once, checking the whole word
\end{itemize}

\begin{tcolorbox}[breakable,colback=teal!4,colframe=teal!35!black,title={Before you move on}]
What does \pa{ਭੰਨਣਾ} mean? (\textquotedblleft{}To smash\textquotedblright{}.) Say it once more.
\end{tcolorbox}
\section[\texorpdfstring{\pa{ਗਹਿਣਾ} (gahiṇā)}{ਗਹਿਣਾ (gahiṇā)}]{\pa{ਗਹਿਣਾ} (gahiṇā) — a piece of jewellery}

\label{lesson:PA-C315-gahina}



\begin{quote}
Before the new one: say the Punjabi for an attack, then the Punjabi for to smash.
\end{quote}

\subsection*{You'll want to know: \pa{ਗਹਿਣਾ}}
\textbf{\pa{ਗਹਿਣਾ}} — \emph{gahiṇā} — \textquotedblleft{}a piece of jewellery\textquotedblright{}.

\begin{quote}
\pa{ਮਾਤਾ} \pa{ਜੀ} \pa{ਦੇ} \pa{ਗਹਿਣੇ} \pa{ਅਲਮਾਰੀ} \pa{ਵਿੱਚ} \pa{ਸਨ।}
\end{quote}

\begin{quote}
\emph{mātā jī de gahiṇe almārī vich san.}
\end{quote}

\begin{quote}
Mother's jewellery was in the cupboard.
\end{quote}

A masculine noun.

\begin{cousinweb}
\textbf{\pa{ਗਹਿਣਾ}} (\emph{gahiṇā}) is inherited from Sanskrit \emph{grahaṇa}, \textquotedblleft{}seizing, holding\textquotedblright{}; CDIAL lists Punjabi \emph{gahiṇā}, \textquotedblleft{}a pledge, a pawn\textquotedblright{}, under it; the sense of a jewel is given there only for Hindi \emph{gahnā} (CDIAL 4364).
\end{cousinweb}

\begin{grammarlens}[title={What you've built}]
One more word for saying what happened.
\end{grammarlens}

\subsection*{Your turn}
\begin{itemize}
\raggedright
  \item \emph{Say it:} \emph{gahiṇā}
  \item \emph{Say it:} \emph{gahiṇā}, once more
  \item \emph{Say it:} \emph{mātā jī de gahiṇe almārī vich san.} — Mother's jewellery was in the cupboard.
  \item \emph{Recall:} say the Punjabi for an attack, then the Punjabi for to smash, then say \emph{gahiṇā} again
  \item \emph{Write it:} copy \textbf{\pa{ਗਹਿਣਾ}} once, checking the whole word
\end{itemize}

\begin{tcolorbox}[breakable,colback=teal!4,colframe=teal!35!black,title={Before you move on}]
What does \pa{ਗਹਿਣਾ} mean? (\textquotedblleft{}A piece of jewellery\textquotedblright{}.) Say it once more.
\end{tcolorbox}
\section[\texorpdfstring{\pa{ਜਾਨ} (jān)}{ਜਾਨ (jān)}]{\pa{ਜਾਨ} (jān) — a life}

\label{lesson:PA-C315-jan}



\begin{quote}
Before the new one: say the Punjabi for to smash, then the Punjabi for a piece of jewellery.
\end{quote}

\subsection*{You'll want to know: \pa{ਜਾਨ}}
\textbf{\pa{ਜਾਨ}} — \emph{jān} — \textquotedblleft{}a life\textquotedblright{}.

\begin{quote}
\pa{ਹਮਲੇ} \pa{ਵਿੱਚ} \pa{ਸਾਡੀ} \pa{ਜਾਨ} \pa{ਬਚੀ।}
\end{quote}

\begin{quote}
\emph{hamle vich sāḍī jān bachī.}
\end{quote}

\begin{quote}
We came through the attack alive.
\end{quote}

A feminine noun.

\begin{cousinweb}
\textbf{\pa{ਜਾਨ}} (\emph{jān}) is a loan, usually traced to Persian \emph{jān}, \textquotedblleft{}life, soul\textquotedblright{}.
\end{cousinweb}

\begin{grammarlens}[title={What you've built}]
One more word for saying what happened.
\end{grammarlens}

\subsection*{Your turn}
\begin{itemize}
\raggedright
  \item \emph{Say it:} \emph{jān}
  \item \emph{Say it:} \emph{jān}, once more
  \item \emph{Say it:} \emph{hamle vich sāḍī jān bachī.} — We came through the attack alive.
  \item \emph{Recall:} say the Punjabi for to smash, then the Punjabi for a piece of jewellery, then say \emph{jān} again
  \item \emph{Write it:} copy \textbf{\pa{ਜਾਨ}} once, checking the whole word
\end{itemize}

\begin{tcolorbox}[breakable,colback=teal!4,colframe=teal!35!black,title={Before you move on}]
What does \pa{ਜਾਨ} mean? (\textquotedblleft{}A life\textquotedblright{}.) Say it once more.
\end{tcolorbox}
\section[(dialogue)]{Looking back: from News from Home to Bandits on the Road}

\label{lesson:PA-R315-review-309-315}



\begin{quote}
Say the Punjabi for a piece of news and an accident.
\end{quote}

\subsection*{Your turn}
Say these aloud:

\begin{itemize}
\raggedright
  \item a piece of news, an accident, a rumour, a din, to die
  \item to slip, a stumble, to get broken, to break something, to get spilt
  \item to seek, a thief, to run away, to snatch, damage
  \item a meeting with someone, to invite, to knock on a door, a host, a gift
  \item a wedding, a groom's wedding party, a bridegroom, to hand out, congratulations
  \item to be born, a small child, to bring up, a son's son, a daughter's son
  \item a bandit, an attack, to smash, a piece of jewellery, a life
\end{itemize}

\begin{tcolorbox}[breakable,colback=teal!4,colframe=teal!35!black,title={Before you move on}]
A piece of news? (\textbf{\pa{ਖ਼ਬਰ}}, \emph{khabar}.) An accident? (\textbf{\pa{ਹਾਦਸਾ}}, \emph{hādsā}.) A rumour? (\textbf{\pa{ਅਫ਼ਵਾਹ}}, \emph{afvāh}.) A din? (\textbf{\pa{ਰੌਲਾ}}, \emph{raulā}.) To die? (\textbf{\pa{ਮਰਨਾ}}, \emph{marnā}.) To slip? (\textbf{\pa{ਫਿਸਲਣਾ}}, \emph{phisalṇā}.) A stumble? (\textbf{\pa{ਠੋਕਰ}}, \emph{ṭhokar}.) To get broken? (\textbf{\pa{ਟੁੱਟਣਾ}}, \emph{ṭuṭṭṇā}.) To break something? (\textbf{\pa{ਤੋੜਨਾ}}, \emph{toṛnā}.) To get spilt? (\textbf{\pa{ਡੁੱਲ੍ਹਣਾ}}, \emph{ḍullhṇā}.) To seek? (\textbf{\pa{ਭਾਲਣਾ}}, \emph{bhālṇā}.) A thief? (\textbf{\pa{ਚੋਰ}}, \emph{chor}.) To run away? (\textbf{\pa{ਨੱਸਣਾ}}, \emph{nassṇā}.) To snatch? (\textbf{\pa{ਖੋਹਣਾ}}, \emph{khohṇā}.) Damage? (\textbf{\pa{ਨੁਕਸਾਨ}}, \emph{nuksān}.) A meeting with someone? (\textbf{\pa{ਮੁਲਾਕਾਤ}}, \emph{mulākāt}.) To invite? (\textbf{\pa{ਸੱਦਣਾ}}, \emph{saddṇā}.) To knock on a door? (\textbf{\pa{ਖੜਕਾਉਣਾ}}, \emph{khaṛkāuṇā}.) A host? (\textbf{\pa{ਮੇਜ਼ਬਾਨ}}, \emph{mezbān}.) A gift? (\textbf{\pa{ਤੋਹਫ਼ਾ}}, \emph{tohfā}.) A wedding? (\textbf{\pa{ਵਿਆਹ}}, \emph{viāh}.) A groom's wedding party? (\textbf{\pa{ਬਰਾਤ}}, \emph{barāt}.) A bridegroom? (\textbf{\pa{ਲਾੜਾ}}, \emph{lāṛā}.) To hand out? (\textbf{\pa{ਵੰਡਣਾ}}, \emph{vaṇḍṇā}.) Congratulations? (\textbf{\pa{ਵਧਾਈ}}, \emph{vadhāī}.) To be born? (\textbf{\pa{ਜੰਮਣਾ}}, \emph{jammṇā}.) A small child? (\textbf{\pa{ਨਿਆਣਾ}}, \emph{niāṇā}.) To bring up? (\textbf{\pa{ਪਾਲਣਾ}}, \emph{pālṇā}.) A son's son? (\textbf{\pa{ਪੋਤਾ}}, \emph{potā}.) A daughter's son? (\textbf{\pa{ਦੋਹਤਾ}}, \emph{dohtā}.) A bandit? (\textbf{\pa{ਡਾਕੂ}}, \emph{ḍākū}.) An attack? (\textbf{\pa{ਹਮਲਾ}}, \emph{hamlā}.) To smash? (\textbf{\pa{ਭੰਨਣਾ}}, \emph{bhannṇā}.) A piece of jewellery? (\textbf{\pa{ਗਹਿਣਾ}}, \emph{gahiṇā}.) A life? (\textbf{\pa{ਜਾਨ}}, \emph{jān}.)
\end{tcolorbox}
